% ============================================================================
\clearpage
\section{Arquitectura y operación diaria}\label{sec:arquitectura}
% ============================================================================

\ficha{Snapshots, parámetros versionados, posiciones y caja, calendario de mercado.}
{Ejecutar la cadena de módulos en horarios fijos, registrar cada resultado y exponerlo por API e
interfaz.}
{Informe diario, recomendación o abstención, órdenes simuladas idempotentes y reconciliadas.}
{Corrida reproducible con el mismo snapshot; ningún fallo aguas arriba se convierte en orden.}

\subsection{Componentes}\label{sec:componentes}

La implementación propuesta separa un motor cuantitativo en Python, almacenamiento de snapshots en
Parquet consultable mediante DuckDB, metadatos de corridas y posiciones en SQLite, una API independiente y
una interfaz web para distribuciones y decisiones (\figref{fig:arquitectura}). Es una propuesta de
arquitectura, sin selección final de versiones ni benchmark de rendimiento.

\begin{figure}[H]
\centering
\resizebox{\linewidth}{!}{% Arquitectura: módulos separados, almacenamiento, registro, API e interfaz.
\begin{tikzpicture}[
  mod/.style={caja=#1,text width=2.05cm,minimum height=10mm},
  bd/.style={cylinder,shape border rotate=90,aspect=0.22,draw=tinta2,fill=grisclaro,line width=0.6pt,
    font=\sffamily\scriptsize,text=tinta,align=center,text width=2.3cm,minimum height=15mm,inner sep=2pt},
  x=1cm,y=1cm]
% Módulos del motor (serpiente)
\node[mod=qazul] (m1) at (3.9,0.9) {Ingestión};
\node[mod=qazul] (m2) at (6.45,0.9) {Control de\\ calidad};
\node[mod=qazul] (m3) at (9.0,0.9) {Superficies};
\node[mod=qazul] (m4) at (11.55,0.9) {Distribuciones};
\node[mod=qazul] (m5) at (14.1,0.9) {Transporte\\ y factores};
\node[mod=qnaranja] (m6) at (14.1,-0.85) {Predicción $\P$};
\node[mod=qnaranja] (m7) at (11.55,-0.85) {Calibración};
\node[mod=qaqua] (m8) at (9.0,-0.85) {Escenarios};
\node[mod=qaqua] (m9) at (6.45,-0.85) {Decisión};
\node[mod=qaqua] (m10) at (3.9,-0.85) {Simulación};
\foreach \a/\b in {m1/m2,m2/m3,m3/m4,m4/m5,m6/m7,m7/m8,m8/m9,m9/m10} \draw[flecha] (\a) -- (\b);
\draw[flecha] (m5) -- (m6);
\node[nota,anchor=south] at (9.0,1.5) {motor cuantitativo en Python: cada módulo con entradas y salidas versionadas};
% Almacenamiento
\node[bd] (s1) at (0.6,0.9) {Snapshots\\ inmutables\\ \tiny Parquet + DuckDB};
\node[bd] (s2) at (0.6,-1.6) {Corridas y\\ posiciones\\ \tiny SQLite};
\draw[flecha] (s1) -- (m1);
% Registro transversal
\node[caja=tenue,text width=12.1cm,minimum height=8mm] (reg) at (9.0,-2.3) {\textbf{Registro}: hashes de entrada,
  versión de modelo, parámetros, hora de corte y motivo de abstención de cada resultado};
\draw[flecha] (reg.west) -- (s2.east |- reg.west);
% Servicios
\node[caja=tenue,text width=2.6cm] (api) at (5.3,-3.75) {API independiente};
\node[caja=tenue,text width=2.6cm] (ui) at (9.0,-3.75) {Interfaz web};
\node[caja=qaqua,text width=2.6cm] (paper) at (12.7,-3.75) {Conector de paper};
\draw[flecha] (reg.south -| api.north) -- (api.north);
\draw[flecha] (api) -- (ui);
\draw[flecha] (reg.south -| paper.north) -- (paper.north);
\end{tikzpicture}
}
\caption{Módulos del motor y servicios. Los colores siguen las tres capas: medición $\Q$ (azul),
predicción $\P$ (naranja) y acción (verde).}
\label{fig:arquitectura}
\end{figure}

\begin{table}[H]
\centering\small
\caption{Responsabilidad de cada módulo. El paquete de referencia \texttt{quantileflow/} del repositorio
contiene versiones mínimas de los bloques matemáticos, usadas para las figuras.}
\label{tab:modulos}
\begin{tabularx}{\linewidth}{@{}l Y l@{}}
\toprule
\encab{Módulo} & \encab{Responsabilidad} & \encab{Referencia}\\
\midrule
Ingestión & Capturar snapshots inmutables con ambos sellos de tiempo & ---\\
Control de calidad & Filtros, estados de confianza y motivos de exclusión & ---\\
Superficies & SSVI restringido, ajuste convexo, desamericanización & \texttt{superficies}, \texttt{opciones}\\
Distribuciones & Breeden--Litzenberger, controles, cuantiles y soporte & \texttt{distribuciones}\\
Transporte y factores & $W_1$, $W_2$, cambios firmados, PCA funcional, filtros & \texttt{transporte}, \texttt{filtrado}\\
Predicción $\P$ & Modelos por objetivo y horizonte & \texttt{pronostico}\\
Calibración & Reglas de puntuación, PIT, conformal adaptativo & \texttt{calibracion}\\
Escenarios y decisión & Escenarios conjuntos, CVaR, costos, región casi óptima & \texttt{decision}\\
Regímenes y evidencia & BOCPD, Sharpe deflactado, bootstrap, particiones & \texttt{regimenes}, \texttt{evidencia}\\
Simulación y registro & Paper, reconciliación y trazabilidad & ---\\
\bottomrule
\end{tabularx}
\end{table}

\subsection{Trazabilidad y reproducibilidad}\label{sec:trazabilidad}

Cada resultado conserva hashes de entrada, versión del modelo, parámetros, hora de corte y motivo de
abstención. Con el mismo snapshot, la corrida se reproduce exactamente. El siguiente registro es un ejemplo
ilustrativo de su forma, no una salida real:

\begin{lstlisting}[style=registro]
{
  "corrida": "2026-09-25T09:45:00-04:00/SPY",
  "hora_corte": "2026-09-25T09:45:00-04:00",
  "snapshot": {"sha256": "3f9a...c21e", "fuente": "proveedor-a", "retraso_s": 0},
  "codigo": {"version": "0.1.0", "commit": "a1b2c3d"},
  "modelos": {"superficie": "ssvi-potencia", "p_h5": "qr-l1", "intervalos": "aci-gamma-0.01"},
  "parametros_sha256": "9d04...7b10",
  "confianza": {"datos": "apto", "ajuste": "degradado", "prediccion": "apto", "decision": "apto"},
  "salida": {"accion": "mantener", "region_peso": [0.30, 0.46], "cvar95_5s": 0.031},
  "abstencion": {"incrementos": "bloqueados", "motivo": "ajuste degradado: bandas de cola anchas"}
}
\end{lstlisting}

\subsection{Calendario diario}\label{sec:calendario}

La captura empieza en la apertura regular; la primera evaluación se hace alrededor de las 09:45 ET si la
calidad alcanza el umbral, con una comprobación posterior a las 10:00 y un registro próximo al cierre
(\figref{fig:calendario}). Las medias sesiones y los cambios de horario ajustan la agenda. Los canales
cierre--apertura y apertura--apertura se mantienen separados.

\begin{figure}[H]
\centering
\resizebox{\linewidth}{!}{% Calendario diario (hora del Este) y canales de cambio.
\begin{tikzpicture}[x=1cm,y=1cm,
  hito/.style={circle,fill=#1,minimum size=2.6mm,inner sep=0pt},
  nota/.style={font=\sffamily\scriptsize,text=tinta2,align=center}]
% Agenda vertical
\draw[guia,line width=0.8pt] (0,0.3) -- (0,-4.6);
\node[hito=qazul] at (0,0) {};
\node[titulo,anchor=east] at (-0.3,0) {09:30};
\node[nota,anchor=west,text=tinta,align=left] at (0.35,0) {\textbf{Apertura regular}: empieza la captura de snapshots};
\node[hito=qazul] at (0,-1.1) {};
\node[titulo,anchor=east] at (-0.3,-1.1) {09:45};
\node[nota,anchor=west,text=tinta,align=left] at (0.35,-1.1) {\textbf{Primera evaluación}, solo si la calidad\\ alcanza el umbral};
\node[hito=qazul] at (0,-2.2) {};
\node[titulo,anchor=east] at (-0.3,-2.2) {10:00};
\node[nota,anchor=west,text=tinta,align=left] at (0.35,-2.2) {\textbf{Comprobación posterior}: ¿la recomendación\\ es estable respecto de la de 09:45?};
\node[hito=qaqua] at (0,-3.3) {};
\node[titulo,anchor=east] at (-0.3,-3.3) {$\approx$15:50};
\node[nota,anchor=west,text=tinta,align=left] at (0.35,-3.3) {\textbf{Registro} próximo al cierre};
\node[hito=tenue] at (0,-4.35) {};
\node[titulo,anchor=east] at (-0.3,-4.35) {16:00};
\node[nota,anchor=west,text=tinta,align=left] at (0.35,-4.35) {Cierre; medias sesiones y cambios de horario\\ ajustan esta agenda};
% Canales de cambio
\begin{scope}[xshift=8.3cm,yshift=-0.4cm]
\node[titulo,anchor=west] at (0,1.0) {Dos canales que no se mezclan};
\fill[grisclaro] (0,-0.9) rectangle (2.2,-0.5);
\fill[grisclaro] (3.2,-0.9) rectangle (5.4,-0.5);
\node[nota,anchor=north] at (1.1,-0.95) {sesión $t-1$};
\node[nota,anchor=north] at (4.3,-0.95) {sesión $t$};
\draw[flecha,draw=qnaranja,line width=0.8pt] (2.2,-0.45) .. controls (2.5,0.15) and (2.9,0.15) .. (3.2,-0.45);
\node[nota,anchor=south] at (2.7,0.2) {cierre $\to$ apertura};
\draw[flecha,draw=qazul,line width=0.8pt] (0.25,-1.35) .. controls (1.4,-2.0) and (2.3,-2.0) .. (3.45,-1.35);
\node[nota,anchor=north] at (1.85,-1.95) {apertura $\to$ apertura};
\node[nota,anchor=north west,text width=6.2cm,align=left] at (0,-2.45) {Con una fuente retrasada 15 minutos, el
  instante \emph{ejecutable} de decisión se mueve en la misma cantidad: la evaluación de 09:45 solo puede
  usar datos hasta 09:30.};
\end{scope}
\end{tikzpicture}
}
\caption{Agenda de una sesión y canales de cambio. Con una fuente retrasada, el instante ejecutable de
decisión se mueve.}
\label{fig:calendario}
\end{figure}

\subsection{Pantalla}\label{sec:pantalla}

La pantalla muestra distribuciones con bandas, el mapa temporal de cuantiles, los cambios de cola y sus
innovaciones, los pronósticos $\P$ diferenciados de $\Q$, la posición actual frente a la propuesta con costos
y riesgo conjunto, la razón de la decisión y el historial de calibración (\figref{fig:pantalla}).

\begin{figure}[H]
\centering
\resizebox{0.95\linewidth}{!}{% Maqueta de la pantalla principal (solo estructura; sin datos).
\begin{tikzpicture}[x=1cm,y=1cm,
  panel/.style={draw=eje,fill=white,rounded corners=2pt,line width=0.6pt,anchor=north west,
    font=\sffamily\scriptsize,text=tinta,align=left,inner sep=5pt},
  etiq/.style={font=\sffamily\scriptsize\bfseries,text=tinta,anchor=north west}]
\draw[eje,fill=grisfondo,rounded corners=3pt,line width=0.7pt] (0,0) rectangle (15.4,-9.2);
% Barra superior
\fill[qazul!10] (0.05,-0.05) rectangle (15.35,-0.75);
\node[etiq] at (0.25,-0.18) {SPY \textperiodcentered{} 25 sep 2026 \textperiodcentered{} corte 09:45 ET};
\node[font=\sffamily\scriptsize,text=tinta,anchor=north east] at (15.2,-0.18)
  {\tikz[baseline=-0.6ex]\node[circle,fill=bueno,minimum size=2mm,inner sep=0pt]{}; datos: apto
   \quad \tikz[baseline=-0.6ex]\node[circle,fill=alerta,minimum size=2mm,inner sep=0pt]{}; ajuste: degradado};
% Fila 1
\node[panel,minimum width=7.45cm,minimum height=3.05cm,text width=7.1cm] at (0.25,-0.95)
  {\textbf{Distribución $\Q$ con bandas}\\ densidad y cuantiles a 30 días; colas no identificadas marcadas};
\node[panel,minimum width=7.45cm,minimum height=3.05cm,text width=7.1cm] at (7.9,-0.95)
  {\textbf{Mapa temporal de cuantiles}\\ últimas 60 sesiones; choques y datos faltantes señalados};
% Fila 2
\node[panel,minimum width=7.45cm,minimum height=2.2cm,text width=7.1cm] at (0.25,-4.2)
  {\textbf{Cambios de cola e innovaciones}\\ $\Delta x_t(u)$ firmado; innovación propia y relativa al mercado};
\node[panel,minimum width=7.45cm,minimum height=2.2cm,text width=7.1cm] at (7.9,-4.2)
  {\textbf{Pronósticos $\P$ (distintos de $\Q$)}\\ cuantiles a 1, 5 y 20 sesiones; calibración reciente};
% Fila 3
\node[panel,minimum width=4.85cm,minimum height=2.45cm,text width=4.5cm] at (0.25,-6.6)
  {\textbf{Posición actual frente a propuesta}\\ región de pesos estables; costos; riesgo conjunto};
\node[panel,minimum width=4.85cm,minimum height=2.45cm,text width=4.5cm] at (5.3,-6.6)
  {\textbf{Razón de la decisión}\\ texto generado a partir de cifras estructuradas; motivo de abstención};
\node[panel,minimum width=4.85cm,minimum height=2.45cm,text width=4.5cm] at (10.35,-6.6)
  {\textbf{Historial de calibración}\\ cobertura, PIT y aciertos por horizonte ya maduro};
\end{tikzpicture}
}
\caption{Maqueta de la pantalla principal (estructura, sin datos). Los estados de confianza aparecen en la
barra superior con icono y etiqueta.}
\label{fig:pantalla}
\end{figure}

\begin{noconcluir}
Un componente de lenguaje puede explicar resultados estructurados, pero no inventar cifras ni decidir
tamaños por texto libre. Todo número que aparezca en la explicación debe provenir del registro de la
corrida.
\end{noconcluir}

\subsection{Conector de paper y órdenes}\label{sec:ordenes}

El conector de paper es un componente separado: órdenes idempotentes, reconciliación de posiciones y caja,
cancelación y estado de ejecuciones (\figref{fig:ordenes}). Inicialmente el producto solo recomienda y
registra; la simulación de ejecución se incorpora después de comprobar el motor.

\begin{figure}[H]
\centering
\resizebox{\linewidth}{!}{% Ciclo de una orden en el conector de paper: validación, idempotencia y reconciliación.
\begin{tikzpicture}[
  est/.style={caja=#1,text width=2.35cm,minimum height=10mm},
  node distance=7mm and 8mm]
\node[est=qaqua] (rec) {Recomendación};
\node[diamond,aspect=2.4,draw=qazul,fill=qazul!7,line width=0.6pt,font=\sffamily\scriptsize,align=center,
  text=tinta,inner sep=1pt,text width=2.4cm,right=of rec] (puerta) {Puerta de\\ validación};
\node[est=tenue,right=of puerta] (prop) {Orden propuesta\\ \tiny clave idempotente};
\node[est=tenue,right=of prop] (env) {Enviada};
\node[est=bueno,below=of env] (ejec) {Ejecutada};
\node[est=tenue,left=of ejec] (parc) {Parcialmente\\ ejecutada};
\node[est=alerta,right=of ejec] (canc) {Cancelada o\\ rechazada};
\node[est=qaqua,below=of parc] (recon) {Reconciliación\\ posiciones y caja};
\node[est=tenue,left=of recon] (regi) {Registro};
\node[est=critico,below=of puerta] (sin) {Sin orden:\\ abstención registrada};
\draw[flecha] (rec) -- (puerta);
\draw[flecha] (puerta) -- node[nota,above]{pasa} (prop);
\draw[flecha] (puerta) -- node[nota,right,align=left]{datos o superficie\\ inválidos, límites} (sin);
\draw[flecha] (prop) -- (env);
\draw[flecha] (env) -- (parc);
\draw[flecha] (env) -- (ejec);
\draw[flecha] (env) -- (canc);
\draw[flecha] (parc) -- (ejec);
\draw[flecha] (parc) -- (recon);
\draw[flecha] (ejec) |- (recon);
\draw[flecha] (canc) |- (recon);
\draw[flecha] (recon) -- (regi);
\draw[flecha] (sin) |- (regi);
\draw[flecha,draw=tenue,densely dashed] (env.north) .. controls +(0,0.9) and +(0,0.9) .. node[nota,above]{reintento con la misma clave: no duplica} (prop.north);
\end{tikzpicture}
}
\caption{Ciclo de una orden simulada. La puerta de validación impide que un fallo de superficie o de datos
se convierta silenciosamente en una orden; la clave idempotente evita duplicados en los reintentos.}
\label{fig:ordenes}
\end{figure}

\begin{criterio}
Operación reproducible en paper durante un periodo acordado, con fallos inyectados (datos faltantes,
proveedor caído, ejecuciones parciales) manejados sin órdenes duplicadas ni posiciones descuadradas.
\end{criterio}
